\section{Conclusion}
\label{sec:conclusion}

We described a reconstruction pipeline whose supervision is a compact, analytic, frozen 2D
Gaussian field per view instead of the source photographs, and we split its evaluation into
two firewalled claims.  Part~I --- the systems claim --- argues that after a one-time
per-view fit, standard 3DGS optimization never needs the images again: supervision becomes a
stream of point queries against kilobyte-scale fields, and the dataset-scale RGB working set
disappears from the reconstruction process.  The mechanism is implemented end-to-end with a
verified point-rendering parity anchor, an unbiased sampled objective, and a hard structural
no-image boundary; \toshow{the claim itself awaits its controlled experiment --- held-out
parity plus measured peak-VRAM reduction under the preregistered protocol of
\cref{sec:vram-protocol}}.  Part~II --- Beam Fusion --- shows that the same captures support
a closed-form tomographic 3D initialization with a principled fusion rule, an observability
theorem fixing the minimum view count, complete contributor lineage, and a compact-only
carrier maturation whose every stage is justified by paired ablation; \toshow{whether the
initialization is \emph{worth having} against the matched no-Beam control is an open,
preregistered question, and a negative answer will be reported as such without touching
Part~I}.

The broader bet is that capture representations --- not model compression and not optimizer
machinery --- are the under-explored axis of scaling 3D reconstruction: once images are
replaced by queryable analytic fields, memory decouples from view count, supervision
decouples from resolution, and the same construction extends to time.  If the confirmatory
program supports Part~I, dome-scale and video-scale captures stop being a memory problem;
if it refutes it, the negative result will delimit exactly how much information per-view
2D Gaussian fits preserve --- either way, the decision rule was written down first.
